\chapter{Aanbevelingen}

\subsection{Brandveiligheidseisen} 
Voor de brandveiligheid moet aan een aantal dingen voldaan worden om in compliantie te zijn met de (toekomstige) brandveiligheidseisen voor batterijen op schepen. Omdat er nog geen internationale regelgeving is voor brandveiligheid bij batterijen op schepen, worden de richtlijnen die de Deense Maritieme Autoriteit heeft opgesteld, gevolgd. Deze worden gezien als de basis voor de toekomstige internationale regelgeving \cite{DMA_battery_operation_ships}. Echter, omdat deze regelgeving nog niet bestaat, is het mogelijk om met goed verstand een aantal van deze systemen later te implementeren, om zo de kosten enigszins te verspreiden.

In de richtlijnen staat het volgende:

\begin{itemize}
    \item De batterijen zullen in een aparte ruimte worden geplaatst.
    \item de batterijruimte zal vrij zijn van hittebronnen en beschermd zijn tegen externe hitte door middel van hittebestendige isolatie.
    \item de batterijruimte zal een actieve ventilatie hebben, 
    \item een geschikt, automatisch brandblussysteem zal aanwezig zijn, welke voldoet aan de door de batterijfabrikant gegeven specificaties.
\end{itemize}